\documentclass[11pt,a4paper]{article}

\usepackage[T1]{fontenc}
\usepackage[utf8]{inputenc}
\usepackage{lmodern}
\input{glyphtounicode}
\pdfgentounicode=1
\usepackage{amsmath,amssymb,amsthm}
\usepackage{booktabs}
\usepackage{microtype}
\usepackage[hidelinks]{hyperref}
\usepackage{doi}

\bibliographystyle{plain}

\newtheorem{theorem}{Theorem}[section]
\newtheorem{lemma}[theorem]{Lemma}
\newtheorem{proposition}[theorem]{Proposition}
\newtheorem{corollary}[theorem]{Corollary}
\newtheorem{remark}[theorem]{Remark}
\newtheorem{definition}[theorem]{Definition}

\newcommand{\Heis}{\mathrm{Heis}_3}
\newcommand{\Sym}{\mathfrak{S}}

\title{A Finite Countermodel to Heisenberg Carrier Selection\\
from Admissibility Axioms}

\author{J\'er\^ome Beau\\
\small Independent Researcher\\
\small \texttt{jerome.beau@cosmochrony.org}\\
\small ORCID: \href{https://orcid.org/0009-0001-7697-7868}{0009-0001-7697-7868}}

\date{2026}

\begin{document}

\maketitle

\begin{abstract}
Two companion papers in the Cosmochrony programme claim that admissibility axioms select the finite
Heisenberg group and its Weil carrier.
Their load-bearing algebraic implication is that a finite faithful irreducible unitary carrier, equipped
with a minimally generating pair exchanged by an involution and having a non-trivial commutator, must have
a central commutator and hence be of Heisenberg type.
We disprove that implication by an explicit six-element countermodel.
Let \(G=\Sym_3\), \(X=(12)\), \(Y=(23)\), and let \(\sigma\) be conjugation by \((13)\).
Then \(\sigma\) is involutive and exchanges \(X\) and \(Y\); the pair minimally generates \(G\); and the
standard action on the two-dimensional zero-sum subspace of \(\mathbb{C}^3\) is faithful, unitary, and
irreducible.
Nevertheless, \([X,Y]\) is a non-trivial three-cycle and is not central, since \(Z(\Sym_3)=\{1\}\).
We also show independently that the published prime-order step reverses the content of Schur's lemma:
a central subgroup acts through one character on an irreducible complex representation and therefore does
not yield the asserted non-trivial eigenspace decomposition.
The algebraic consequences extracted from the admissibility inputs in the published proof do not select
class-two nilpotence, a finite Heisenberg group, or a Weil carrier.
\textbf{Interpretive status.}
The countermodel changes the status of the carrier-selection edge from theorem to postulate; it does not
invalidate results proved after a Heisenberg carrier is supplied.
\end{abstract}

\noindent\textbf{Keywords:}
finite countermodel; carrier selection; Heisenberg group; admissibility; irreducible representation;
Schur's lemma; nilpotence; epistemic status

\clearpage
\tableofcontents

\clearpage
\section{Question and Audit Contract}
\label{sec:contract}

The Foundation paper claims that Axioms A1--A3 together with Born--Infeld parity force
\[
G_n \simeq \Heis(\mathbb{Z}/q\mathbb{Z}),
\qquad
F_n \simeq V_\rho
\]
for a prime \(q\), and describes the identification as a theorem rather than a
remark~\cite{Beau2026Foundation}.
The companion Heisenberg Structure paper supplies the detailed argument and states the same conclusion as
its main theorem~\cite{Beau2026HeisenbergStructure}.

The present note asks only whether the algebraic properties actually available before the centrality step
force that conclusion.
It does not assess the later mathematics carried out after a Heisenberg group and a Weil representation
have been supplied.

\subsection{Properties used before centrality}

The published argument extracts the following contract from the axioms:

\begin{enumerate}
  \item a finite-dimensional complex carrier \(F\);
  \item a faithful irreducible unitary action of a finite group \(G\) on \(F\);
  \item a generating pair \(X,\sigma(X)\), with no proper subpair generating \(G\);
  \item an involutive automorphism \(\sigma\) exchanging the two generators;
  \item a non-trivial commutator \([X,\sigma(X)]\).
\end{enumerate}

The next step asserts that the commutator is central because any non-trivial higher commutator would be
``excess structure'' incompatible with minimality~\cite{Beau2026HeisenbergStructure}.
But a group generated by two elements need not be nilpotent of class two.
The countermodel below makes the distinction exact.

\subsection{Convention}

Permutations are composed from right to left.
For group elements \(a,b\), the commutator is
\[
[a,b]=aba^{-1}b^{-1}.
\]

\section{The Six-Element Countermodel}
\label{sec:countermodel}

\begin{theorem}[Finite countermodel]
\label{thm:countermodel}
Let
\[
G=\Sym_3,\qquad X=(12),\qquad Y=(23),\qquad \tau=(13),
\]
and define
\[
\sigma(g)=\tau g\tau^{-1}.
\]
Let \(G\) act by coordinate permutation on
\[
W=\bigl\{(z_1,z_2,z_3)\in\mathbb{C}^3:
z_1+z_2+z_3=0\bigr\}.
\]
Then:
\begin{enumerate}
  \item \(\sigma\) is an involutive automorphism and \(\sigma(X)=Y\), \(\sigma(Y)=X\);
  \item \(X\) and \(Y\) form a minimal generating pair for \(G\);
  \item the representation of \(G\) on \(W\) is two-dimensional, faithful, unitary, and irreducible;
  \item \([X,Y]\neq 1\);
  \item \([X,Y]\notin Z(G)\).
\end{enumerate}
Consequently, the five properties listed in Section~\ref{sec:contract} do not force class-two nilpotence or
a Heisenberg group.
\end{theorem}

\begin{proof}
Because \(\tau^2=1\), conjugation by \(\tau\) satisfies \(\sigma^2=\mathrm{id}\).
Conjugation relabels the entries of a transposition, hence
\[
\sigma((12))=(\tau(1)\,\tau(2))=(32)=(23)=Y
\]
and similarly \(\sigma(Y)=X\).

The elements \(X\) and \(Y\) generate all of \(\Sym_3\): their product \(XY=(123)\), so the subgroup they
generate contains a transposition and a three-cycle and therefore has six elements.
Neither singleton \(\{X\}\) nor \(\{Y\}\) generates \(G\), because each generates a subgroup of order two.
The pair is therefore minimal.

The coordinate-permutation action preserves the standard Hermitian product on \(\mathbb{C}^3\), and \(W\)
is invariant under it.
Thus the restricted action on \(W\) is unitary.
Since \(W\) is the kernel of the non-zero linear functional
\((z_1,z_2,z_3)\mapsto z_1+z_2+z_3\), it has dimension two.

To prove faithfulness, suppose \(g\in\Sym_3\) acts trivially on \(W\).
It then fixes both \(e_1-e_2\) and \(e_2-e_3\).
A coordinate permutation fixing these two vectors fixes each of the three coordinate positions, so \(g=1\).

To prove irreducibility, suppose that \(W\) has a \(G\)-invariant complex line \(L\).
The restrictions of the involutions \(X\) and \(Y\) to \(L\) are scalars in \(\{1,-1\}\).
Because \(X\) and \(Y\) are conjugate, these scalars are equal.
If both act as \(+1\), every vector of \(L\) is fixed by \(G\), hence has three equal coordinates; the
zero-sum condition then makes it zero.
If both act as \(-1\), \(L\) lies in the intersection of the \(-1\)-eigenspaces
\[
\mathbb{C}(e_1-e_2)\cap\mathbb{C}(e_2-e_3)=\{0\}.
\]
Both cases contradict the existence of a line.
Since \(W\) is two-dimensional, it is irreducible.

Finally, \(X^{-1}=X\) and \(Y^{-1}=Y\), so
\[
[X,Y]=XYXY=(XY)^2=(132)\neq 1.
\]
It is not central: for example,
\[
X(132)X^{-1}=(123)\neq(132).
\]
This completes the verification.
\end{proof}

\begin{corollary}[Failure of the centrality implication]
\label{cor:centrality}
Minimal generation by an involutively conjugate pair, faithful irreducibility, unitarity, and a non-trivial
commutator do not imply that the commutator is central.
\end{corollary}

\begin{proof}
Theorem~\ref{thm:countermodel} satisfies every premise and violates the conclusion.
\end{proof}

\section{Independent Failure of the Prime-Order Step}
\label{sec:schur}

The published proof next argues that a proper non-trivial subgroup of the centre would decompose the
irreducible carrier into several non-zero character eigenspaces, contradicting non-factorisability.
That conclusion is incompatible with Schur's lemma.
For completeness, the required statement is proved directly.

\begin{lemma}[Central subgroups on an irreducible carrier]
\label{lem:central-character}
Let \(V\) be a finite-dimensional irreducible complex representation of a group \(G\), and let
\(H\leq Z(G)\).
Then \(H\) acts on \(V\) through a single character.
Equivalently, in the decomposition
\[
V=\bigoplus_{\chi\in\widehat H}V_\chi,
\]
at most one summand is non-zero.
\end{lemma}

\begin{proof}
For any \(h\in H\), the complex operator representing \(h\) has an eigenvalue \(\lambda_h\).
Its eigenspace is \(G\)-invariant because \(h\) is central, so irreducibility makes that eigenspace all of
\(V\).
Thus every \(h\) acts as the scalar \(\lambda_h\), and \(h\mapsto\lambda_h\) is a single character of \(H\).
\end{proof}

\begin{proposition}[No prime-order conclusion]
\label{prop:no-prime}
The existence of a proper non-trivial subgroup \(H\leq Z(G)\) does not contradict irreducibility of a complex
\(G\)-representation.
Therefore irreducibility alone does not imply that \(Z(G)\), or a specified central cyclic subgroup, has
prime order.
\end{proposition}

\begin{proof}
By Lemma~\ref{lem:central-character}, the whole carrier may transform through one non-trivial character of
\(H\).
The number of characters of \(H\) does not imply that more than one occurs.
The asserted stable decomposition is therefore trivial rather than forbidden.
\end{proof}

\section{Where the Published Classification Imports Its Conclusion}
\label{sec:classification}

The countermodel closes the issue at the centrality step, but two later inputs are also logically prior to
the classification they are used to support.

\subsection{Central-extension data}

The classification step invokes extensions of
\[
(\mathbb{Z}/q\mathbb{Z})^2
\quad\text{by}\quad
\mathbb{Z}/q\mathbb{Z}
\]
and the presentation
\[
X^q=Y^q=Z^q=1,\qquad [X,Y]=Z,\qquad Z\in Z(G).
\]
To apply that classification, one must already know the exponent, the abelianisation, the cyclic centre,
centrality of the commutator, and the non-degeneracy of its pairing.
The audit contract in Section~\ref{sec:contract} supplies none of those data.
Classifying central extensions after they have been specified does not show that the original axioms specify
one.

\subsection{Character parity}

The finite character involution \(c\mapsto q-c\) presupposes a finite cyclic centre of order \(q\) and its
character labelling.
Continuum parity of a Born--Infeld candidate does not by itself construct that finite carrier.
Using the character involution before the carrier has been selected therefore cannot serve as an independent
selection argument.

\subsection{Conditional role of Stone--von Neumann}

Once a finite Heisenberg group, a non-trivial central character, and irreducibility are supplied, the finite
Stone--von Neumann theorem identifies the corresponding carrier.
That is a valid conditional final step.
It does not establish the group presentation, the order of the centre, or the character selection required
to invoke it.

\section{Status Consequences}
\label{sec:status}

\subsection{What is refuted}

Theorem~\ref{thm:countermodel} refutes the implication
\[
\begin{gathered}
\text{faithful irreducible unitary carrier}
+\text{minimal involutive conjugate pair}\\
\text{}+\text{non-trivial commutator}\\
\not\Longrightarrow\\
\text{central commutator and Heisenberg carrier}.
\end{gathered}
\]
Accordingly, the claim that A1--A3 and Born--Infeld parity uniquely select
\(\Heis(\mathbb{Z}/q\mathbb{Z})\) and \(V_\rho\) is not established by the published derivation.
At the present axiomatic level, the Heisenberg/Weil carrier is a supplied realisation.

\subsection{What survives}

The obstruction does not alter:
\begin{itemize}
  \item the elementary fact that an irreducible unitary carrier of dimension greater than one cannot be
  generated by commuting operators;
  \item the classification and representation theory of a finite Heisenberg group once that group is
  supplied;
  \item any theorem in the O-series, Q-series, or companion notes whose statement is explicitly internal to
  a chosen Heisenberg graph or Weil carrier.
\end{itemize}

Those downstream results become conditional on carrier selection only when they are read as consequences of
the admissibility axioms.
Their internal mathematics is unchanged.

\subsection{What would reopen selection}

The \(S_3\) model can be excluded by adding a hypothesis such as class-two nilpotence, a specified central
extension, or an equivalent non-degenerate symplectic commutator law.
But such a hypothesis is precisely the missing selection input.
It would require an independent motivation or derivation rather than being inferred from the absence of
``excess structure''.

\section{Conclusion}
\label{sec:conclusion}

The six-element group \(\Sym_3\) provides a complete finite countermodel to the algebraic selection step used
to derive a Heisenberg carrier from admissibility.
Its standard two-dimensional representation is faithful, unitary, and irreducible; its minimal generating
pair is exchanged by an involution and has a non-trivial commutator; yet that commutator is not central.
The separate prime-order argument also fails because an irreducible carrier contains only one character of
a central subgroup, not one non-zero eigenspace for every available character.

The correct logical status is therefore:
\[
\text{admissibility axioms}
\;\not\Longrightarrow\;
\Heis(\mathbb{Z}/q\mathbb{Z})
\quad\text{under the published argument}.
\]
A Heisenberg/Weil carrier remains a mathematically productive choice, but its unique selection is postulated
until an additional class-two or symplectic input is independently established.

\subsection*{Interpretive outlook}

\textbf{Interpretive status.}
This result is an obstruction to a bridge, not a refoundation of the mathematical programme.
It separates two questions that had been merged: what follows on a Heisenberg carrier, and why that carrier
should describe the admissible fibre.
The first question retains the corpus of internal results already proved.
The second becomes an explicit open selection problem.
Making that edge visible permits exploratory use of the carrier without presenting the exploration as an
axiomatic necessity.

\bibliography{references}

\end{document}
